\section{Introduction}
Evaluating public support for charging infrastructure requires connecting the provision of facilities to the service they sustain. Construction and deployment milestones make the available stock visible to public authorities, whereas drivers experience that investment through access to charging and its continued operation. Expansion can therefore serve a coverage objective before it produces frequent activity, particularly where access is being extended or demand is still developing. Conversely, a falling activity-to-facility ratio may accompany activity growth that is slower than stock growth or an absolute contraction in activity. Those cases have different implications for evaluating support: one concerns the pace at which service follows expansion, while the other concerns whether observed service is being sustained at all. An evaluation that treats both as a single form of underperformance loses this distinction and cannot state clearly what a construction programme has delivered.

The infrastructure and adoption literature provides a first reason for keeping the distinction open. Charging infrastructure and electric-vehicle adoption can reinforce one another through indirect network effects, but the strength and timing of that relationship depend on policy design and market conditions \citep{li2017,springel2021,burra2024}. Evidence on income, charger type and regional context likewise shows that deployment is translated into adoption unevenly \citep{gifford2026,gabehart2026}; Chinese evidence distinguishes the role of vehicle subsidies from infrastructure subsidies in strengthening infrastructure--adoption links \citep{hu2025_charging}. This interdependence makes deployment relevant to market development, while leaving a separate question about the activity recorded by the available network. Vehicle adoption, installed stock and charging activity measure different stages of that relationship. A facility indicator can describe the opportunity made available, while an activity measure describes a recorded service response conditional on what is observed.

Research on charging service adds a second layer of heterogeneity. Utilization differs with charger type, venue, network conditions and the time over which it is measured; occupancy, delivered energy, connection duration and profitability each use a different denominator and answer a different question \citep{borlaug2023,hecht2022,wolbertus2018}. Station-level evidence links use to accessibility, built environment, land use and temporal context, so two cities with similar facility counts need not offer the same service mix \citep{zhao2024_utilization,du2024_accessibility,zheng2025_landuse}. Access and observed use can diverge further when demand is censored by capacity, concentrated by habitual routines or affected by reliability \citep{huttel2023,weekx2024,weekx2025_habits,rempel2024,yu2025_reliability}. The relevant comparison is therefore not whether one universal utilization metric is available, but whether the outcome, exposure and time support are aligned. In this study, recorded activity, observed facilities and activity relative to facilities are kept as related but non-equivalent objects.

The unresolved evaluation problem is consequently one of scale and time as well as measurement. Network-wide activity intensity is a facility-weighted average of city-specific intensities, whereas an equal-city summary gives each city the same weight. A majority of cities can improve while network-wide intensity declines if the improvements occur in small networks or are offset by larger changes elsewhere. Similarly, a category assigned from one endpoint window says little about whether the same city sustains that pattern later. Fixed-member comparisons separate within-city change from changing coverage; recorded-state comparisons retain source-coded zeros without treating missing observations as zeros; and paired categories reveal whether a relative-growth pattern persists or changes. These distinctions allow a more precise question than whether the network is simply ``underused'': how does recorded service develop relative to the observed stock across cities and horizons, and which cities or recording states account for the network result?

That question also has an institutional side. Public support can specify reporting or platform connection, operation and maintenance, continued operation, or payment eligibility, and these duties concern different parts of the post-construction service chain. Policy-mix research emphasizes that instruments and their interactions, rather than document counts alone, shape what a programme asks of its recipients \citep{rogge2016}. Research on Chinese charging subsidies and collaborative governance similarly points to operating performance, feedback and administrative coordination as policy concerns \citep{hu2026,zhang2024,guo2026}; systems-based analysis highlights data sharing and enforceable management arrangements across public and private actors \citep{tilly2025}. Textual requirements do not establish compliance or a causal policy effect. They do, however, identify responsibilities that can be checked against observed activity, equipment condition or continuity, making the evaluation question institutionally concrete.

We address these distinctions by linking 2016--2023 city-week records from one anonymized operator network with locally applicable policy documents. A consistent measurement of activity and observed facilities allows city-level changes to be compared with facility-weighted network intensity. Holding city membership fixed separates growth from widening archive coverage, while retaining source-coded-zero records reveals how much facility scale lies outside continuously positive activity. Pairing cities across periods then distinguishes a common growth pattern from a persistent one. Conditional annual and quarterly relationships address a further question: whether activity is associated with facility scale in the same proportion, and whether the finite-horizon association exceeds its contemporaneous component. The policy material brings these comparisons back to the responsibilities specified for supported infrastructure, distinguishing what texts require from the outcome associations that can be observed.

The contribution is an empirical account of why geographic breadth, network scale and temporal persistence lead to different evaluations of charging expansion. In the 2022--2023 comparison, 117 of 225 cities have source-coded-zero activity at both endpoints, yet they represent only 3.47\% of endpoint facilities. Counting these cities equally therefore gives a very different description from using facility-weighted network intensity. Growth types are also changeable: 17 of the 27 early faster-growth cities in the paired comparison subsequently shift to slower positive growth or contraction. These descriptive results complement PPML and quarterly conditional estimates of proportionality; the annual log-OLS interval still permits unit elasticity, and the additional quarterly follow-through interval includes zero. Read alongside the responsibilities prescribed in local policies, the findings support separate evaluation of coverage provision, the pace and continuity of recorded activity, and the operating conditions attached to public support.
